\section*{Conclusion}
\addcontentsline{toc}{section}{Conclusion}

Ce rapport a permis d'analyser la politique RSE d'Iveco Group et d'identifier les points d'ancrage du numérique responsable au sein du site de Haute-Goulaine.

Sur le plan de la gouvernance, Iveco Group dispose d'un dispositif structuré : priorités RSE claires, instances dédiées (Conseil d'administration, comité ESG, SLT, équipes Energy et EEHS), double matérialité alignée sur les ESRS et une déclaration consolidée via ESGeo. Les reconnaissances ESG (S\&P Global CSA top 5\,\%, EcoVadis Gold, CDP A/A-, ISS ESG B-) valident cette maturité. Toutefois, le score CDP Forêts (C) et la prédominance des émissions Scope 3 dans l'empreinte carbone rappellent que la décarbonation dépendra avant tout de la transition énergétique des véhicules vendus et de l'engagement des fournisseurs.

Sur le volet social, la politique DEI est pilotée activement : objectifs chiffrés (30\,\% de femmes dans les fonctions de bureau d'ici 2028), certification EDGE, suivi du handicap et de la diversité culturelle. La santé et la sécurité affichent des résultats très positifs, avec une réduction de 51\,\% du taux de fréquence de blessures par rapport à 2019. L'engagement communautaire, mesuré par le cadre B4SI, démontre une capacité à transformer les contributions financières et en temps en impacts mesurables auprès de près de 150\,000 bénéficiaires.

Enfin, le numérique responsable apparaît comme le maillon encore partiellement adressé. La politique de confidentialité intègre un solide \textit{Privacy by Design}, et le groupe utilise déjà des outils numériques (IMDS, ESGeo, Web Academy) pour réduire ses impacts. En revanche, aucun programme Green IT spécifique à l'informatique interne n'est documenté dans les ressources. Les recommandations formulées pour l'équipe de Haute-Goulaine (éco-conception logicielle, cycle de vie du matériel, achats responsables, maîtrise de l'énergie et des données, sensibilisation et volontariat) proposent des leviers concrets, mesurables et réalistes.

La prochaine étape consisterait à valider ces recommandations avec la direction informatique groupe, à établir une base de référence énergétique et matérielle au premier trimestre 2026, puis à intégrer les indicateurs proposés dans le tableau de bord du site. Cette démarche permettrait d'aligner l'informatique locale sur les ambitions RSE d'Iveco Group, tout en renforçant la valeur ajoutée d'un ingénieur en informatique face aux enjeux de durabilité.
